% Generated draft: requires mathematical and source review.
\documentclass[11pt]{article}
\usepackage[margin=25mm]{geometry}
\usepackage{fontspec}
\setmainfont{Latin Modern Roman}
\usepackage{amsmath,amssymb,mathtools,unicode-math}
\setmathfont{Latin Modern Math}
\usepackage{enumitem,needspace,xurl,hyperref}
\setlength{\parindent}{0pt}
\setlength{\parskip}{5pt}
\setlength{\emergencystretch}{3em}
\newcommand{\N}{\mathbb{N}}
\newcommand{\Z}{\mathbb{Z}}
\newcommand{\Q}{\mathbb{Q}}
\newcommand{\R}{\mathbb{R}}
\newcommand{\C}{\mathbb{C}}
\newcommand{\F}{\mathbb{F}}
\newcommand{\A}{\mathbb{A}}
\newcommand{\PP}{\mathbb{P}}
\newcommand{\E}{\mathbb{E}}
\newcommand{\eps}{\varepsilon}
\newcommand{\dd}{\,\mathrm{d}}
\begin{document}
\section*{UnsolvedMath ID: 30004766}
\textbf{Source problem code:} catalogue entry
\subsection{Definitions and mathematical statement}
For a field \(k\) and a \(k\)-algebra \(A\) with bilinear multiplication, its derivation Lie algebra is
\[
\operatorname{Der}_k(A)=\{D\in\operatorname{End}_k(A):D(ab)=D(a)b+aD(b)\text{ for all }a,b\in A\},
\]
with bracket \([D,E]=D\circ E-E\circ D\). This definition explains the usual algebra interpretation of derivations; the supplied record does not specify the objects whose derivations are intended, so it does not establish that this interpretation exhausts the problem's scope. A Lie algebra is called simple here if it is nonzero, nonabelian, and has no nonzero proper ideals. The intended meaning of ``non-Chevalley'' is not defined in the supplied record and must be fixed before selecting the input class.

For a represented Lie algebra \(L\subseteq\mathfrak{gl}_k(V)=\operatorname{End}_k(V)\) and a specified subgroup \(G\leq\operatorname{GL}_k(V)\), the group normalizer is
\[
N_G(L)=\{g\in G:gLg^{-1}=L\}.
\]
A different possible target is the Lie algebra normalizer
\[
N_{\mathfrak{gl}_k(V)}(L)=\{X\in\mathfrak{gl}_k(V):[X,L]\subseteq L\}.
\]
These definitions describe distinct possible readings; the source record does not select an ambient group, an ambient Lie algebra, or a representation.

\textbf{Source question.} Can normalizers of the non-Chevalley simple Lie algebras that arise as derivation algebras be constructed efficiently?

Schematically, the requested construction should apply to every member of the intended class of non-Chevalley simple derivation Lie algebras. However, that class, its ground fields and characteristics, its representations, the ambient normalizer, and the computational input and output conventions remain unspecified. In particular, ``efficiently'' cannot yet be replaced by a precise running-time bound. The question is therefore retained as an unresolved algorithmic formulation, rather than strengthened to a polynomial-time assertion or restricted to an invented domain.
\subsection{Short English statement}
The problem asks for an efficient method to construct normalizers of non-Chevalley simple Lie algebras occurring as derivation algebras. A precise algorithmic statement requires the intended fields, derivation objects, representations, ambient normalizers, and meaning of efficiency. The supplied registry labels the problem open; this draft does not independently verify that status.
\subsection{Sources}
\begin{itemize}
\item[1.] Supplied local UnsolvedMath registry record, permanent ID \(30004766\), title ``Efficient Normalizers of Non-Chevalley Derivation Lie Algebras,'' exact source problem code \texttt{catalogue entry}, in the collection ``Oberwolfach Reports Open Problems.'' The record gives proposed year \(2021\) and names no individual proposer.
\item[2.] Supplied report citation: \emph{Computational Group Theory (2022)}, Oberwolfach Reports, DOI \texttt{10.4171/owr/2021/38}, pp.\ 2027--2087. \url{https://doi.org/10.4171/owr/2021/38}. The literature-source label is ``Oberwolfach Report 38/2021.'' The supplied citation's \(2022\) designation and the report's \(2021\) designation are retained without reconciliation.
\item[3.] Supplied literature source: \emph{Tensor Isomorphism by Derivations and Densors}. \url{https://arxiv.org/abs/2005.04046}. The supplied record does not identify authors or a specific version.
\item[4.] Literature assessment dated exactly \texttt{2026-08-22}: the supplied curation labels the problem open and reports that the cited work provides efficient normalizer constructions for important Chevalley-type cases while leaving the requested construction for non-Chevalley simple Lie algebras arising as derivation algebras as a future problem. This is supplied dated curation, not independently verified literature assessment. The background explicitly requests independent verification of time-sensitive status and cited claims.
\item[5.] The supplied registry assigns statement status \texttt{corrected\_verified} and reports that the official Oberwolfach report was checked directly and that the clean statement preserves the complete mathematical target while restoring definitions and removing extraction artifacts or source-dependent wording. These are supplied verification claims; no independent source consultation is asserted here.
\end{itemize}
\end{document}
